\section*{Речник на основните термини}
\addcontentsline{toc}{section}{Речник на основните термини}
\label{sec:rechnik}

{\small
\begin{longtable}{@{}T{4.1cm}L{12.45cm}@{}}
Електронно изявление & Словесно изявление, представено в цифрова форма чрез общоприет
стандарт за преобразуване, разчитане и представяне на информацията; може да съдържа и
несловесна информация (чл.~2 ЗЕДЕУУ).\\
Електронен документ & Всяко съдържание, съхранявано в електронна форма, по-специално
текстов или звуков, визуален или аудио-визуален запис (чл.~3, т.~35 eIDAS; чл.~3, ал.~1
ЗЕДЕУУ).\\
Автор на електронното изявление & Физическото лице, което в изявлението се сочи като негов
извършител (чл.~4 ЗЕДЕУУ).\\
Титуляр на електронното изявление & Лицето, от името на което е извършено изявлението;
може да бъде и юридическо лице (чл.~4 ЗЕДЕУУ).\\
Титуляр на електронен подпис & Физическото лице, което създава електронен подпис (чл.~3,
т.~9 eIDAS).\\
Електронен подпис & Данни в електронна форма, които се добавят към други данни в електронна
форма или са логически свързани с тях и които титулярят използва, за да се подписва
(чл.~3, т.~10 eIDAS).\\
Усъвършенстван електронен подпис & Електронен подпис, който е свързан по уникален начин с
титуляря, може да го идентифицира, създаден е с данни под негов контрол и позволява да
бъде открита всяка последваща промяна в подписаните данни (чл.~3, т.~11 и чл.~26 eIDAS).\\
Квалифициран електронен подпис (КЕП) & Усъвършенстван електронен подпис, създаден от
устройство за създаване на квалифициран електронен подпис и основан на квалифицирано
удостоверение за електронни подписи (чл.~3, т.~12 eIDAS).\\
Данни за създаване на електронен подпис & Уникални данни, които титулярят използва за
създаването на подписа; при КЕП това е частният ключ (чл.~3, т.~13 eIDAS).\\
Квалифицирано удостоверение за електронен подпис & Електронен атестат, който свързва
данните за валидиране на подписа с физическо лице, издава се от квалифициран доставчик и
отговаря на приложение~I на eIDAS (чл.~3, т.~14 и~15 eIDAS).\\
Устройство за създаване на КЕП & Устройство, което отговаря на изискванията на
приложение~II на eIDAS, например смарт карта или сървърно устройство на доставчика
(чл.~3, т.~23 eIDAS).\\
Квалифициран доставчик на удостоверителни услуги & Доставчик, който предоставя една или
повече квалифицирани удостоверителни услуги и е получил квалифицирания си статут от
надзорен орган, в България~-- от КРС (чл.~3, т.~20 eIDAS; чл.~32 ЗЕДЕУУ).\\
Валидиране & Процесът на проверка и потвърждаване на валидността на данните в електронна
форма в съответствие с eIDAS (чл.~3, т.~41 и чл.~32 eIDAS).\\
Квалифициран електронен времеви печат & Данни в електронна форма, които свързват други
данни с конкретен момент и се ползват от презумпция за точност на датата и часа и за
цялост на данните (чл.~3, т.~33 и~34 и чл.~41, пар.~2 eIDAS).\\
Хеш стойност & Кратка контролна сума с фиксирана дължина, която се изчислява от
съдържанието на документа; всяка промяна в документа води до различна хеш стойност.\\
Формална доказателствена сила & Качеството на подписания частен документ да доказва, че
изявленията в него са направени от лицето, което го е подписало (чл.~180 ГПК).\\
Оспорване на истинността & Процесуално действие, с което страната твърди, че документът не
изхожда от лицето, посочено като негов автор (чл.~193 ГПК; вж. чл.~93, т.~6 НК).\\
\end{longtable}
}
% longtable увеличава брояча на таблиците и без надпис; връщаме го, за да започне номерацията от 1.
\addtocounter{table}{-1}
